\chapter{Normativa y accesibilidad}

\section{Protección de datos}

El proyecto se orienta a una aplicación potencial en rehabilitación cognitiva, por lo que debe contemplar desde el diseño la protección de datos personales y, especialmente, de datos de salud. Durante el desarrollo descrito en esta memoria no se ha trabajado con datos reales de pacientes. Los registros analizados corresponden a vehículos autónomos simulados, métricas de entrenamiento y resultados técnicos.

En una futura validación clínica, cualquier uso con personas usuarias deberá ajustarse al Reglamento General de Protección de Datos y a la Ley Orgánica 3/2018 de Protección de Datos Personales y garantía de los derechos digitales~\cite{rgpd,lopdgdd}. Será necesario definir responsable del tratamiento, finalidad, base jurídica, minimización de datos, conservación, consentimiento informado cuando proceda, medidas de seudonimización y control de acceso. También deberá diferenciarse entre datos necesarios para la evaluación terapéutica y métricas técnicas internas del simulador.

\section{Inteligencia artificial, ética y confidencialidad}

El sistema desarrollado incorpora técnicas de inteligencia artificial, pero en el estado actual trabaja con vehículos simulados y no toma decisiones clínicas sobre personas. Si en el futuro se integra en un entorno sanitario o de rehabilitación con usuarios reales, deberá revisarse su encaje dentro del Reglamento Europeo de Inteligencia Artificial, especialmente por su posible relación con salud, seguridad, supervisión humana, trazabilidad y gestión del riesgo~\cite{aiact}. En ese escenario no bastaría con que el algoritmo funcionara técnicamente: sería necesario documentar finalidad, límites de uso, métricas de rendimiento, intervención profesional y procedimientos de parada o corrección.

Antes de cualquier estudio con pacientes o personas usuarias sería recomendable la revisión por un comité de ética o por el órgano que corresponda dentro de la entidad responsable. También debería definirse un protocolo de anonimización o seudonimización, separación entre identificadores personales y métricas del simulador, periodos de conservación, control de accesos, copias de seguridad y registro de cambios. La memoria actual no describe datos clínicos reales, por lo que estas medidas se formulan como requisitos para fases posteriores.

El repositorio y los datos de desarrollo no se publican como material abierto en esta memoria. La razón no es ocultar resultados, sino proteger assets, Blueprints, modelos, trazas de colaboración y posibles elementos sujetos a licencias o acuerdos de trabajo con entidades externas. Para defensa académica se aporta trazabilidad mediante anexos, tablas, figuras, referencias, CSV analizados y descripción técnica, evitando enlaces públicos a repositorios o ficheros que puedan exponer información privada o recursos redistribuibles.

\section{Seguridad y uso responsable}

El simulador no sustituye por sí mismo una evaluación clínica ni una autorización administrativa para conducir. Su función debe entenderse como apoyo a profesionales, entrenamiento controlado y herramienta de observación. Cualquier conclusión sobre capacidad real de conducción debería depender de especialistas y procedimientos establecidos.

Desde el punto de vista técnico, el sistema autónomo genera tráfico artificial. Por tanto, sus fallos no suponen riesgo físico directo durante el entrenamiento virtual, pero sí pueden afectar a la validez de la experiencia. Es importante que las sesiones con pacientes utilicen modelos suficientemente probados para evitar comportamientos incoherentes que confundan al usuario o introduzcan estímulos no deseados.

\section{Licencias de assets externos}

El proyecto utiliza una mesh de semáforo descargada de Fab, identificada en la documentación del trabajo mediante la URL \href{https://fab.com/s/c00389c4b9f2}{\texttt{https://fab.com/s/c00389c4b9f2}}. Fab distingue entre licencias Creative Commons Attribution para determinados assets gratuitos y licencia Standard para assets gratuitos o de pago~\cite{fablicenses}. En la licencia Standard existen los niveles Personal y Professional, que dependen de los ingresos brutos de actividad comercial en los últimos 12 meses, pero ambos conceden el mismo alcance de uso sobre el asset~\cite{fabstandard}.

Según el resumen oficial de la Fab Standard License, el asset puede utilizarse de forma privada o comercial, modificarse para incorporarlo a un proyecto, utilizarse con herramientas compatibles no limitadas a Unreal Engine y distribuirse comercialmente integrado dentro del proyecto final~\cite{fabstandard}. También puede compartirse con colaboradores que trabajen en el mismo proyecto, por ejemplo mediante un repositorio privado. La restricción principal es que no puede revenderse ni redistribuirse como asset independiente, ni gratis ni de pago, ni facilitar que terceros lo hagan~\cite{fabstandard}.

Para esta memoria, el uso realizado es compatible con esa finalidad: la mesh actúa como elemento visual integrado en el simulador y no como producto autónomo. Aunque Fab indica que no es obligatorio acreditar al publicador bajo la licencia Standard, se conserva la referencia al origen por trazabilidad académica y para facilitar una revisión posterior. Si en el futuro el simulador se distribuye fuera del entorno académico, deberá comprobarse de nuevo la ficha concreta del asset descargado, conservar la licencia aceptada en el momento de adquisición y evitar incluir el archivo fuente de la mesh en repositorios públicos donde pueda extraerse como recurso independiente.

\clearpage
\section{Accesibilidad}

La accesibilidad del simulador debe analizarse en relación con las personas a las que se dirige, que pueden presentar alteraciones de atención, memoria, percepción o movilidad. En este TFG no se ha rediseñado la interfaz utilizada por el paciente, por lo que no puede afirmarse que el producto completo haya sido validado como accesible. La aportación realizada se concentra en el tráfico autónomo y afecta de forma indirecta a la carga cognitiva de cada sesión.

No consta en las fuentes revisadas una certificación formal de accesibilidad ni una auditoría completa según una norma específica. Las WCAG 2.2 ofrecen principios útiles de perceptibilidad, operabilidad, comprensibilidad y robustez para interfaces digitales~\cite{wcag22}, pero su aplicación a un simulador 3D no web debe interpretarse con prudencia. En este proyecto pueden servir como referencia parcial para menús, instrucciones y paneles futuros, no como evidencia de cumplimiento global.

El número de vehículos, los puntos del mapa y la frecuencia de situaciones normativas deben poder ajustarse para que la dificultad aumente de forma progresiva. También es importante repetir un mismo escenario bajo condiciones comparables: si el tráfico cambia sin control, resulta difícil saber si una variación en el resultado procede del paciente o de una sesión más exigente. La separación entre entrenamiento y test de la IA y el registro de eventos sientan una base técnica para construir esas repeticiones, aunque todavía falta trasladar estos parámetros a una configuración manejable por el profesional.

Desde la perspectiva de accesibilidad, el tráfico autónomo debe comportarse como una variable graduable. Un perfil de baja dificultad debería reducir densidad, velocidad y número de decisiones simultáneas; un perfil avanzado podría introducir cruces ocupados, esperas en ceda, rotondas y vehículos próximos. Esta graduación no depende solo del modelo neuronal, sino también del gestor de spawns, del modo de ejecución y de los indicadores que se presenten al especialista. Por ello, la accesibilidad futura debe contemplar tanto la interfaz del paciente como las herramientas de configuración clínica.

La interacción debe contemplar periféricos adaptados, instrucciones breves y legibles y una presentación que no añada estímulos irrelevantes. El rendimiento gráfico y la duración de la prueba también deben vigilarse para reducir fatiga o mareo. Por último, los datos exportados no deberían mostrarse al especialista como un volcado técnico de CSV, sino como indicadores comprensibles y vinculados a maniobras concretas. Estos requisitos deberán revisarse con profesionales y usuarios de ASPAYM antes de utilizar el simulador en una validación clínica.
